problem:  
Let \((M, \mathcal{D}, g)\) be a compact, equiregular **contact sub-Riemannian manifold** with horizontal bundle \(\mathcal{D}\) and volume measure \(\mu\) induced by the contact form. Denote by \(\mathrm{Ric}_{\mathrm{ABR}}(\lambda)\) the **Agrachev–Barilari–Rizzi canonical Ricci curvature** evaluated along a normal sub-Riemannian geodesic with initial covector \(\lambda \in T^*M\).  

Assume that there exists a constant \(K>0\) such that for every unit covector \(\lambda\) generating a normal geodesic,  
\[
\mathrm{Ric}_{\mathrm{ABR}}(\lambda(t)) \ge K, \quad \forall\, t \text{ along the geodesic.}
\]  
Moreover, suppose that the metric-measure space \((M, d_{\mathrm{subR}}, \mu)\) satisfies the **Measure Contraction Property** \(MCP(K,N)\) for some finite \(N>0\), where \(d_{\mathrm{subR}}\) is the sub-Riemannian distance.  

**Question:**  
Under these conditions, must \(M\) be **contact-metrically equivalent** (up to scaling of the horizontal metric) to a **compact Sasakian space form of positive curvature**, i.e., locally sub-Riemannian isometric to a quotient of the standard sphere \(S^{2n+1}\) with its canonical contact distribution?  
Equivalently, does the coexistence of a positive ABR curvature bound and a strong synthetic curvature bound \(MCP(K,N)\) enforce rigidity of the sub-Riemannian geometry to the Sasakian model?  

Why is it a "good" problem:  
This problem achieves precise formulation while remaining open and pivotal to several directions in modern geometric analysis.  

- **Mathematical precision:** It uses the established Agrachev–Barilari–Rizzi notion of canonical Ricci curvature (defined via Jacobi curves along normal extremals) and a standard optimal-transport curvature condition \(MCP(K,N)\), ensuring that all hypotheses are well-defined and operational.  
- **Depth and originality:** The conjecture asks whether intrinsic positive curvature (ABR Ricci) aligned with synthetic measure convexity \(MCP(K,N)\) is strong enough to enforce **rigidity to the Sasakian model**, serving as a sub-Riemannian analogue of the **Bonnet–Myers–Obata rigidity theorem** known in Riemannian geometry.  
- **Interdisciplinary bridge:** It connects sub-Riemannian geometric control theory (via ABR curvature), analytic optimal transport (via MCP), and global differential geometry (via Sasakian classification), forging an unprecedented synthesis of these fields.  
- **Profound implications:** A resolution would yield a unified understanding of curvature bounds and geometric rigidity in nonholonomic spaces, potentially impacting optimal transport, metric geometry, and geometric PDEs on sub-Riemannian manifolds.  
- **Extreme simplicity of statement:** The question is expressed in a few lines of geometric language familiar across different subfields, yet its resolution would require deep integration of analysis, geometry, and topology—making it a textbook example of “narrow entrance, profound depth.”